We proposed a novel approach to estimating the number of new genetic variants that will be observed in a follow-up sequencing study of two potentially heterogeneous populations.

Several opportunities remain for future work. First, we rely on numerical integration for our predictions. Even in the two-population case, these computations can be somewhat expensive. 
\internal{We observe in \cref{app:extension}} that the form of our rate measure in \cref{eq:proposed_prior} is straightforward to extend to more populations. However, 
we expect further computational (and statistical) challenges when considering higher numbers of populations. %We suspect that alternative approaches may prove fruitful.
In addition, we explored only one possible option for generalizing the beta process to higher dimensions. Although our experiments validated our approach, even better performance might be possible with an alternative model -- especially one carefully aligned with biological prior knowledge. \revision{For instance, one drawback of our current rate measure (\cref{eq:proposed_prior}) is its evident asymmetry across populations. It remains to be seen if a symmetric rate measure might be designed to have at least as many desirable properties as \cref{eq:proposed_prior} enjoys (including conjugacy, \cref{req:finite_mean,req:inf_mass}, power law growth rates, empirical performance).}
